\begin{derivation}[有限作用时间的一量子发射概率]
对矩形作用窗完成时间积分并取模平方，就得到正文\eqnrefs{eq:universal_boson_finite_time_emission} 的 $\sinc^2$ 谱核。 在相互作用绘景中，电子算符与模式算符分别按自由哈密顿量 演化。 对电子本征态，

\begin{align*}
{}_V\!\langle f_e|\delta\hat h_\lambda^\dagger(t)|i_e\rangle_V
&=\ee^{\ii(E_f-E_i)t/\hbar}
{}_V\!\langle f_e|\delta\hat h_\lambda^\dagger|i_e\rangle_V,\\
\langle n_\lambda+1|\hat a_\lambda^\dagger(t)|n_\lambda\rangle
&=\ee^{+\ii\omega_\lambda t}\sqrt{n_\lambda+1}.
\end{align*}
因此发射末态 $|f_e;\mathbf n+\mathbf e_\lambda\rangle_V$ 与初态之间的瞬时矩阵元为
\begin{align*}
\langle f_e;\mathbf n+\mathbf e_\lambda|\hat H_I(t)|i_e;\mathbf n\rangle
&=M_{fi,\lambda}^{\rm em}\sqrt{n_\lambda+1}
\exp\!\left\{\ii\left[\frac{E_f-E_i}{\hbar}+\omega_\lambda\right]t\right\}\\
&=M_{fi,\lambda}^{\rm em}\sqrt{n_\lambda+1}
\ee^{\ii\Delta_{fi}^{\rm em}t}.
\end{align*}
Dyson 级数的一阶项于是给
\[
C_f^{(1)}(T)
=-\frac{\ii}{\hbar}M_{fi,\lambda}^{\rm em}\sqrt{n_\lambda+1}
\int_0^T\dd t\,\ee^{\ii\Delta_{fi}^{\rm em}t}.
\]
时间积分可以写成
\[
\int_0^T\dd t\,\ee^{\ii\Delta t}
=T\ee^{\ii\Delta T/2}\sinc(\Delta T/2),
\]
故
\[
C_f^{(1)}(T)
=-\ii\mathcal G_{fi,\lambda}^{\rm em}\sqrt{n_\lambda+1}
T\ee^{\ii\Delta_{fi}^{\rm em}T/2}
\sinc(\Delta_{fi}^{\rm em}T/2).
\]
取模平方时中间相位 $\exp(\ii\Delta T/2)$ 消失，得到\eqnrefs{eq:universal_boson_finite_time_emission}。其中 $n_\lambda+1$ 直接来自 $|\langle n_\lambda+1|a_\lambda^\dagger|n_\lambda\rangle|^2$；即使 $n_\lambda=0$，产生算符仍有非零矩阵元。吸收支路的末态为 $|f_e;\mathbf n-\mathbf e_\lambda\rangle_V$。此时
\begin{align*}
{}_V\!\langle f_e|\delta\hat h_\lambda(t)|i_e\rangle_V
&=\ee^{\ii(E_f-E_i)t/\hbar}M_{fi,\lambda}^{\rm abs},\\
\langle n_\lambda-1|\hat a_\lambda(t)|n_\lambda\rangle
&=\ee^{-\ii\omega_\lambda t}\sqrt{n_\lambda},
\end{align*}
所以总相位的角频率失谐为
\[
 \Delta_{fi}^{\rm abs}
 =\frac{E_f-E_i}{\hbar}-\omega_\lambda
 =\frac{E_f-E_i-\hbar\omega_\lambda}{\hbar}.
\]
因此
\begin{align*}
 C_{f,\rm abs}^{(1)}(T)
 &=-\ii\mathcal G_{fi,\lambda}^{\rm abs}\sqrt{n_\lambda}
 T\,\ee^{\ii\Delta_{fi}^{\rm abs}T/2}
 \sinc(\Delta_{fi}^{\rm abs}T/2),\\
 P_{i\to f}^{\rm abs}(T)
 &=n_\lambda|\mathcal G_{fi,\lambda}^{\rm abs}|^2T^2
 \sinc^2(\Delta_{fi}^{\rm abs}T/2).
\end{align*}
这与发射支路的 $n_\lambda+1$ 和 $\Delta_{fi}^{\rm em}$ 构成两条独立的有限时间结果。
\end{derivation}

\begin{derivation}[两次量子发射的时间排序与中间态]
取时间无关耦合在有限散射窗内绝热开关，并写发射顶角的相互作用绘景矩阵元
\[
\langle m,n+1|\hat V_+(t_1)|i,n\rangle
=\sqrt{n+1}\,V^{\rm em}_{mi}
\exp\!\left[
\frac{\ii(E_m-E_i+\hbar\omega_0)t_1}{\hbar}
\right],
\]
\[
\langle f,n+2|\hat V_+(t_2)|m,n+1\rangle
=\sqrt{n+2}\,V^{\rm em}_{fm}
\exp\!\left[
\frac{\ii(E_f-E_m+\hbar\omega_0)t_2}{\hbar}
\right].
\]
代入二阶 Dyson 振幅，定义
\[
\Omega_1=\frac{E_m-E_i+\hbar\omega_0}{\hbar},
\qquad
\Omega_2=\frac{E_f-E_m+\hbar\omega_0}{\hbar},
\]
则时间核为
\[
I_T(\Omega_2,\Omega_1)
=\int\dd t_2\int^{t_2}\dd t_1\,
\ee^{\ii\Omega_2t_2}\ee^{\ii\Omega_1t_1}.
\]
时间排序区域 $t_1<t_2$ 最适合改用“总时间”和“两次顶角的时间间隔”作坐标：
\begin{equation*}
 t\equiv t_2,
 \qquad
 \tau\equiv t_2-t_1\ge0,
 \qquad
 t_1=t-\tau,
 \qquad
 \dd t_2\,\dd t_1=\dd t\,\dd\tau.
\end{equation*}
指数因子于是精确分解成
\begin{align*}
 \ee^{\ii\Omega_2t_2}\ee^{\ii\Omega_1t_1}
 &=\ee^{\ii\Omega_2t}
   \ee^{\ii\Omega_1(t-\tau)}\\
 &=\ee^{\ii(\Omega_1+\Omega_2)t}
   \ee^{-\ii\Omega_1\tau}.
\end{align*}
$t$ 积分只检查整个过程的总能量守恒，而 $\tau\ge0$ 记录第一顶角之后中间态能够传播的持续时间；两个积分因此分别编码整体守恒律与中间态传播。

为把半轴积分定义成分布，在相对时间上加入收敛因子 $\ee^{-\epsilon\tau}$，$\epsilon>0$，最后取 $\epsilon\to0^+$：
\begin{align*}
 I_\epsilon(\Omega_2,\Omega_1)
 &=\int_{-\infty}^{\infty}\dd t\,
   \ee^{\ii(\Omega_1+\Omega_2)t}
   \int_0^\infty\dd\tau\,
   \ee^{-(\epsilon+\ii\Omega_1)\tau}\\
 &=2\pi\delta(\Omega_1+\Omega_2)
   \frac1{\epsilon+\ii\Omega_1}\\
 &=2\pi\delta(\Omega_1+\Omega_2)
   \frac{-\ii}{\Omega_1-\ii\epsilon}.
\end{align*}
总时间部分给
\begin{equation*}
2\pi\delta(\Omega_1+\Omega_2)
=2\pi\hbar\,
\delta(E_f-E_i+2\hbar\omega_0).
\end{equation*}
另一方面，二阶 Dyson 前因子是 $(-\ii/\hbar)^2=-1/\hbar^2$。因此相对时间积分与 Dyson 前因子的乘积为
\begin{align*}
 -\frac1{\hbar^2}
 \frac{-\ii}{\Omega_1-\ii0^+}
 &=\frac{\ii}{\hbar^2}
 \frac{\hbar}{E_m-E_i+\hbar\omega_0-\ii0^+}\\
 &=-\frac{\ii}{\hbar}
 \frac1{E_i-E_m-\hbar\omega_0+\ii0^+}.
\end{align*}
于是完整二阶 $S$ 矩阵元具有标准的形式
\begin{align*}
 \langle f,n+2|S^{(2)}|i,n\rangle
 ={}&-2\pi\ii\,
 \delta(E_f-E_i+2\hbar\omega_0)\\
 &\times\sqrt{(n+1)(n+2)}
 \sum_m
 \frac{V^{\rm em}_{fm}V^{\rm em}_{mi}}
 {E_i-E_m-\hbar\omega_0+\ii0^+}.
\end{align*}
去掉散射矩阵中共同的 $-2\pi\ii\delta(E_f-E_i+2\hbar\omega_0)$ 后，约化二阶振幅正是
\[
\mathcal M_{fi}^{(2,\,2{\rm em})}
=\sqrt{(n+1)(n+2)}
\sum_m
\frac{V^{\rm em}_{fm}V^{\rm em}_{mi}}
{E_i-E_m-\hbar\omega_0+\ii0^+}.
\]
若电子中间谱连续，$\sum_m$ 展开为对完整动量、自旋和其它量子数的积分。利用
\[
\frac1{x+\ii0^+}
=\mathcal P\frac1x-\ii\pi\delta(x),
\]
可见中间分母的主值部分对应离壳虚路径，而 $\delta(x)$ 部分表示中间态本身可以在壳传播。后一情形若在有限脉冲带宽内占显著权重，就应恢复显式级联通道，而不应继续把全过程压成局域二阶有效顶角。
\end{derivation}

\begin{derivation}[有限时间谱核与 Lorentz 线形的积分]
保留有限时间的谱核，对 Lorentz 线形直接积分，以求出概率从短时间到长时间的变化。将谱核写成两时间积分，$T^2\sinc^2(\Delta T/2)=\int_0^T\dd t\int_0^T\dd t'\,\ee^{\ii\Delta(t-t')}$，再使用 Lorentz 线形的 Fourier 积分
\begin{equation*}
 \int_{-\infty}^{\infty}\dd\Delta\,
 \frac{J_0\ee^{\ii\Delta u}}{1+\Delta^2\tau_c^2}
 =\frac{\pi J_0}{\tau_c}\ee^{-|u|/\tau_c}.
\end{equation*}
两时间正方形区域中，相同时间差 $u>0$ 对应的线段长度为 $T-u$，故
\begin{align*}
 P^{(2)}(T)&=\frac{2\pi J_0}{\tau_c}
 \int_0^T(T-u)\ee^{-u/\tau_c}\dd u\\
 &=2\pi J_0[T-\tau_c(1-\ee^{-T/\tau_c})].
\end{align*}
因此恒定速率是有限时间积分的渐近结果。谱阈值、窄峰或离散末态改变了被积分的 $\mathcal J_{\rm tr}$ 时，应先重算这个积分，再判断是否存在黄金律区。
\end{derivation}
